% Blueprint: Cayley graph connected iff generators generate the group
% Created by Numina

\begin{document}

\section{Cayley graph connected iff generators generate the group}

Throughout, $G$ is a group (written multiplicatively) and $S \subseteq G$ is a set
of generators. We prove that the (multiplicative) Cayley graph of $G$ with respect
to $S$ is connected if and only if $S$ generates $G$ as a group, i.e.
$\overline{\langle S \rangle} = G$.

\begin{definition}[Multiplicative Cayley graph]
    \label{def:cayley}
    \lean{SimpleGraph.mulCayley}
    \leanfile{LeanEval/Combinatorics/CayleyConnected.lean}
    \leanok
    The \emph{multiplicative Cayley graph} $\mathrm{Cay}(G, S)$ is the simple graph
    \texttt{SimpleGraph.mulCayley}~$S$ on the vertex set $G$: the graph obtained from
    the relation $a \mathrel{R} b \iff \exists g \in S,\ a \cdot g = b$, so that two
    distinct vertices $a, b \in G$ are adjacent precisely when $a \mathrel{R} b$ or
    $b \mathrel{R} a$.
\end{definition}

\begin{lemma}[Adjacency in the Cayley graph]
    \label{lem:adj_iff}
    \lean{LeanEval.Combinatorics.mulCayley_adj_iff}
    \leanfile{LeanEval/Combinatorics/CayleyConnected.lean}
    \uses{def:cayley}
    \leanok
    For $a, b \in G$, the vertices $a$ and $b$ are adjacent in $\mathrm{Cay}(G,S)$
    if and only if $a \neq b$ and there exists $g \in S$ with $a \cdot g = b$ or
    $b \cdot g = a$.
\end{lemma}
\begin{proof}
    \uses{def:cayley}
    Immediate from the definition of the graph associated to a relation
    (\texttt{SimpleGraph.fromRel\_adj}): adjacency in $\mathrm{Cay}(G,S)$ unfolds
    to $a \neq b$ together with the symmetric closure of the defining relation.
\end{proof}

\begin{lemma}[Left translation preserves adjacency]
    \label{lem:adj_mul_left}
    \lean{LeanEval.Combinatorics.mulCayley_adj_mul_left}
    \leanfile{LeanEval/Combinatorics/CayleyConnected.lean}
    \uses{lem:adj_iff}
    \leanok
    For all $h, a, b \in G$, if $a$ and $b$ are adjacent in $\mathrm{Cay}(G,S)$
    then $h \cdot a$ and $h \cdot b$ are adjacent.
\end{lemma}
\begin{proof}
    \uses{lem:adj_iff}
    Using the adjacency characterization, adjacency of $a$ and $b$ gives $a \neq b$
    and some $g \in S$ with $a \cdot g = b$ or $b \cdot g = a$. Left multiplication
    by $h$ is injective, so $h \cdot a \neq h \cdot b$, and from $a \cdot g = b$ we
    get $(h \cdot a) \cdot g = h \cdot b$ (and symmetrically). Hence $h \cdot a$ and
    $h \cdot b$ are adjacent.
\end{proof}

\begin{lemma}[Left translation preserves reachability]
    \label{lem:reachable_mul_left}
    \lean{LeanEval.Combinatorics.mulCayley_reachable_mul_left}
    \leanfile{LeanEval/Combinatorics/CayleyConnected.lean}
    \uses{lem:adj_mul_left}
    \leanok
    For all $h, a, b \in G$, if $b$ is reachable from $a$ in $\mathrm{Cay}(G,S)$
    then $h \cdot b$ is reachable from $h \cdot a$.
\end{lemma}
\begin{proof}
    \uses{lem:adj_mul_left}
    By the previous lemma the map $x \mapsto h \cdot x$ is a graph homomorphism
    $\mathrm{Cay}(G,S) \to \mathrm{Cay}(G,S)$. Graph homomorphisms send walks to
    walks, hence preserve reachability (\texttt{SimpleGraph.Reachable.map}).
\end{proof}

\begin{lemma}[Adjacent vertices differ by a generator]
    \label{lem:adj_mem_closure}
    \lean{LeanEval.Combinatorics.mulCayley_inv_mul_mem_closure_of_adj}
    \leanfile{LeanEval/Combinatorics/CayleyConnected.lean}
    \uses{lem:adj_iff}
    \leanok
    If $a$ and $b$ are adjacent in $\mathrm{Cay}(G,S)$, then $a^{-1} \cdot b \in
    \overline{\langle S \rangle}$, the subgroup generated by $S$.
\end{lemma}
\begin{proof}
    \uses{lem:adj_iff}
    By \texttt{SimpleGraph.mulCayley\_adj}, adjacency of $a$ and $b$ gives directly
    $a^{-1} \cdot b \in S$ or $b^{-1} \cdot a \in S$. In the first case
    $a^{-1} \cdot b \in S \subseteq \overline{\langle S \rangle}$. In the second case
    $a^{-1} \cdot b = (b^{-1} \cdot a)^{-1}$, which lies in
    $\overline{\langle S \rangle}$ since $b^{-1} \cdot a$ does and the closure is
    closed under inverses.
\end{proof}

\begin{lemma}[Reachability implies membership in the closure]
    \label{lem:mem_closure_of_reachable}
    \lean{LeanEval.Combinatorics.mulCayley_inv_mul_mem_closure_of_reachable}
    \leanfile{LeanEval/Combinatorics/CayleyConnected.lean}
    \uses{lem:adj_mem_closure}
    \leanok
    If $b$ is reachable from $a$ in $\mathrm{Cay}(G,S)$, then $a^{-1} \cdot b \in
    \overline{\langle S \rangle}$.
\end{lemma}
\begin{proof}
    \uses{lem:adj_mem_closure}
    Reachability is the reflexive--transitive closure of adjacency
    (\texttt{SimpleGraph.reachable\_iff\_reflTransGen}); induct on it. The reflexive
    base gives $a^{-1} \cdot a = 1 \in \overline{\langle S \rangle}$. For the
    inductive step from $a$ to $b$ followed by an edge $b \sim c$, write
    $a^{-1} \cdot c = (a^{-1} \cdot b)(b^{-1} \cdot c)$; the first factor lies in the
    closure by the induction hypothesis and the second by
    Lemma~\ref{lem:adj_mem_closure}, and the closure is closed under products.
\end{proof}

\begin{lemma}[Generators are reachable from the identity]
    \label{lem:reachable_generator}
    \lean{LeanEval.Combinatorics.mulCayley_reachable_one_generator}
    \leanfile{LeanEval/Combinatorics/CayleyConnected.lean}
    \uses{lem:adj_iff}
    \leanok
    For every $g \in S$, the vertex $g$ is reachable from $1$ in $\mathrm{Cay}(G,S)$.
\end{lemma}
\begin{proof}
    \uses{lem:adj_iff}
    If $g = 1$ this is reflexivity of reachability. Otherwise $1 \neq g$ and
    $1 \cdot g = g$ with $g \in S$, so $1$ and $g$ are adjacent by
    Lemma~\ref{lem:adj_iff}, and a single edge is a walk.
\end{proof}

\begin{lemma}[Reachability from $1$ is closed under products]
    \label{lem:reachable_one_mul}
    \lean{LeanEval.Combinatorics.mulCayley_reachable_one_mul}
    \leanfile{LeanEval/Combinatorics/CayleyConnected.lean}
    \uses{lem:reachable_mul_left}
    \leanok
    For $x, y \in G$, if $x$ and $y$ are each reachable from $1$ in
    $\mathrm{Cay}(G,S)$, then $x \cdot y$ is reachable from $1$.
\end{lemma}
\begin{proof}
    \uses{lem:reachable_mul_left}
    From reachability of $y$ from $1$ and Lemma~\ref{lem:reachable_mul_left}
    (translating by $x$), $x \cdot y$ is reachable from $x \cdot 1 = x$. Since $x$ is
    reachable from $1$, transitivity of reachability gives $x \cdot y$ reachable
    from $1$.
\end{proof}

\begin{lemma}[Reachability from $1$ is closed under inverses]
    \label{lem:reachable_one_inv}
    \lean{LeanEval.Combinatorics.mulCayley_reachable_one_inv}
    \leanfile{LeanEval/Combinatorics/CayleyConnected.lean}
    \uses{lem:reachable_mul_left}
    \leanok
    For $x \in G$, if $x$ is reachable from $1$ in $\mathrm{Cay}(G,S)$, then
    $x^{-1}$ is reachable from $1$.
\end{lemma}
\begin{proof}
    \uses{lem:reachable_mul_left}
    Translating reachability of $x$ from $1$ by $x^{-1}$
    (Lemma~\ref{lem:reachable_mul_left}) shows $x^{-1} \cdot x = 1$ is reachable from
    $x^{-1} \cdot 1 = x^{-1}$. By symmetry of reachability, $x^{-1}$ is reachable
    from $1$.
\end{proof}

\begin{lemma}[Membership in the closure implies reachability]
    \label{lem:reachable_of_mem_closure}
    \lean{LeanEval.Combinatorics.mulCayley_reachable_one_of_mem_closure}
    \leanfile{LeanEval/Combinatorics/CayleyConnected.lean}
    \uses{lem:reachable_generator, lem:reachable_one_mul, lem:reachable_one_inv}
    \leanok
    If $g \in \overline{\langle S \rangle}$, then $g$ is reachable from $1$ in
    $\mathrm{Cay}(G,S)$.
\end{lemma}
\begin{proof}
    \uses{lem:reachable_generator, lem:reachable_one_mul, lem:reachable_one_inv}
    Induct on membership in $\overline{\langle S \rangle}$
    (\texttt{Subgroup.closure\_induction}). Generators are reachable from $1$ by
    Lemma~\ref{lem:reachable_generator}, and $1$ is reachable from itself. The product
    case is Lemma~\ref{lem:reachable_one_mul} and the inverse case is
    Lemma~\ref{lem:reachable_one_inv}, applied to the inductive hypotheses.
\end{proof}

\begin{lemma}[Connectivity as reachability from the identity]
    \label{lem:connected_iff_reachable_one}
    \lean{LeanEval.Combinatorics.mulCayley_connected_iff_reachable_one}
    \leanfile{LeanEval/Combinatorics/CayleyConnected.lean}
    \leanok
    The graph $\mathrm{Cay}(G,S)$ is connected if and only if every vertex $g \in G$
    is reachable from $1$.
\end{lemma}
\begin{proof}
    A graph (on a nonempty vertex type) is connected iff there is a vertex from which
    every vertex is reachable (\texttt{SimpleGraph.connected\_iff\_exists\_forall\_reachable}).
    If every vertex is reachable from $1$, take that vertex to be $1$. Conversely, if
    every vertex is reachable from some $v$, then for any $g$ both $v$ and $g$ are
    reachable from $v$, so by symmetry and transitivity of reachability $g$ is
    reachable from $1$ (which is itself reachable from $v$).
\end{proof}

\begin{theorem}[Cayley graph connected iff generators generate]
    \label{thm:main}
    \lean{LeanEval.Combinatorics.mulCayley_connected_iff_closure_eq_top}
    \leanfile{LeanEval/Combinatorics/CayleyConnected.lean}
    \uses{lem:connected_iff_reachable_one, lem:mem_closure_of_reachable, lem:reachable_of_mem_closure}
    \leanok
    The Cayley graph $\mathrm{Cay}(G,S)$ is connected if and only if
    $\overline{\langle S \rangle} = G$, i.e. $S$ generates $G$.
\end{theorem}
\begin{proof}
    \uses{lem:connected_iff_reachable_one, lem:mem_closure_of_reachable, lem:reachable_of_mem_closure}
    By Lemma~\ref{lem:connected_iff_reachable_one}, connectivity is equivalent to:
    every $g \in G$ is reachable from $1$. By \texttt{Subgroup.eq\_top\_iff'},
    $\overline{\langle S \rangle} = G$ is equivalent to: every $g \in G$ lies in
    $\overline{\langle S \rangle}$. So it suffices to show, for each $g$, that $g$ is
    reachable from $1$ iff $g \in \overline{\langle S \rangle}$. The forward direction
    is Lemma~\ref{lem:mem_closure_of_reachable} with $a = 1$ (so $1^{-1} \cdot g = g$),
    and the reverse direction is Lemma~\ref{lem:reachable_of_mem_closure}.
\end{proof}

\end{document}
